\ifdefined\NoetherPaperTwentyTwoSOneBodyOnly
\else
\documentclass[11pt]{article}
\usepackage[a4paper,margin=1.45cm]{geometry}
\usepackage[T1]{fontenc}
\usepackage[utf8]{inputenc}
\usepackage{lmodern}
\usepackage[french]{babel}
\usepackage{amsmath,amssymb,mathtools}
\usepackage[protrusion=true,expansion=false]{microtype}
\setlength{\parindent}{1.25em}
\setlength{\parskip}{0.10em}
\allowdisplaybreaks
\setlength{\emergencystretch}{12em}
\sloppy\hbadness=10000\vbadness=10000\hfuzz=2pt\vfuzz=10pt\raggedbottom
\makeatletter\renewcommand\footnotesize{\@setfontsize\footnotesize{7.5}{8.7}\selectfont}\makeatother
\providecommand{\Amod}{\mathfrak A}
\providecommand{\Bmod}{\mathfrak B}
\providecommand{\Gmod}{\mathfrak G}
\providecommand{\Mmod}{\mathfrak M}
\providecommand{\Nmod}{\mathfrak N}
\providecommand{\Rmod}{\mathfrak R}
\providecommand{\Smod}{\mathfrak S}
\providecommand{\mideal}{\mathfrak m}
\providecommand{\nideal}{\mathfrak n}
\providecommand{\gideal}{\mathfrak g}
\providecommand{\hideal}{\mathfrak h}
\providecommand{\jideal}{\mathfrak j}
\providecommand{\tideal}{\mathfrak t}
\providecommand{\aideal}{\mathfrak a}
\providecommand{\bideal}{\mathfrak b}
\providecommand{\bb}{\mathfrak b}
\providecommand{\cc}{\mathfrak c}
\providecommand{\zero}[1]{\equiv 0\pmod{#1}}
\providecommand{\Norm}{N}

\begin{document}
\fi

% Unit: Paper 22 opening and \S{}1. Direct French translation from the German final-audited source slice.

\setcounter{footnote}{0}
\section*{22. Sur la théorie des idéaux polynomiaux et des résultantes}
\begin{center}
Math. Ann. 88 (1923), p.~53--79\\[0.9em]
{\Large\bfseries Sur la théorie des idéaux polynomiaux et des résultantes.}\\[1.0em]
Par\\[0.45em]
Kurt Hentzelt \(\dagger\).\\[0.8em]
Édité par Emmy Noether à Göttingen\footnotemark.\\[0.6em]
\rule{4em}{0.4pt}
\end{center}

Il s'agit dans ce qui suit du \emph{problème général de la théorie de l'élimination, déterminer les zéros communs de tous les polynômes d'un idéal de polynômes}; ou, ce qui revient au même, les zéros d'un nombre fini quelconque de polynômes qui forment une base de l'idéal; en même temps apparaîtra une interprétation conceptuelle des multiplicités qui interviennent. Les coefficients des polynômes peuvent être attribués à un corps quelconque; les zéros appartiennent alors au corps algébriquement clos qui en est dérivé. En outre, l'idéal peut être supposé transformé (\S{}~3) sans restriction de généralité. Le résultat essentiel est alors le suivant:

\emph{À tout idéal $\mideal$ de polynômes on peut associer une forme résultante:}
\[
 R_{\mideal}=R^{(1)}(x_1,\ldots,x_n)\cdots
 R^{(i)}(x_i,\ldots,x_n)\cdots R^{(n)}(x_n)\zero{\mideal},
\]
\footnotetext{Kurt Hentzelt est porté disparu depuis octobre 1914 devant Dixmuiden et doit être compté parmi les morts. Le mémoire donné ici est une adaptation entièrement libre de la partie la plus essentielle de sa dissertation \og Zur Theorie der Formenmoduln und Resultanten \fg{}, avec laquelle il obtint son doctorat à l'été 1914 auprès de E. Fischer à Erlangen. Cette dissertation, composée tout entière sur le fondement de ses propres idées, est construite sans lacunes; mais lemme succède à lemme, tous les concepts sont décrits par des formules à quatre et cinq indices, le texte fait presque entièrement défaut, de sorte que les plus grandes difficultés s'opposent à la compréhension. Il n'a lui-même plus eu le temps d'accomplir le remaniement projeté. Je reproduis le travail sous une forme purement conceptuelle; il en résulte une grande simplification des démonstrations, qui remontent partout, dans leurs idées fondamentales, à Hentzelt, et, je l'espère, la beauté du travail devient manifeste. -- Les parties de la dissertation qui, une base étant donnée, concernent la question de la formation, en un nombre fini d'étapes, des fonctions qui apparaissent doivent être réservées à une publication séparée. (E. N.)}
\emph{de telle sorte que $R_{\mideal}$ s'annule pour tous les zéros de $\mideal$, et seulement pour ceux-ci. Si $\nideal$ est un diviseur de $\mideal$, et si les formes résultantes $R_{\nideal}$ et $R_{\mideal}$ coïncident, alors les idéaux $\nideal$ et $\mideal$ coïncident aussi.}

La seconde partie du théorème principal montre que la forme résultante fournit les zéros avec leur multiplicité caractéristique. Cette seconde partie est analogue au fait que deux idéaux d'un corps de nombres algébrique, dont l'un est un diviseur de l'autre, coïncident si et seulement si leurs \emph{normes} coïncident; la raison interne des deux théorèmes est également la même.

En effet, $\mideal$ peut être conçu comme un module de formes linéaires $\mathfrak M_{i-1}$, en considérant chaque polynôme comme une forme linéaire dans les produits de puissances de $x_1,\ldots,x_{i-1}$ et en adjoignant $x_{i+1},\ldots,x_n$ au domaine des coefficients. Le facteur résultant $R^{(i)}$ devient alors égal à la \emph{norme} du module fondamental correspondant (\S{}~1) $\Gmod_{i-1}$ par rapport à $\mathfrak M_{i-1}$; on obtient ainsi aussitôt une \emph{interprétation de la multiplicité -- produit du degré et de l'exposant d'un facteur de $R^{(i)}$ -- par le nombre de classes modulo $\mathfrak M_{i-1}$ linéairement indépendantes}\footnote{Ces derniers résultats montrent leur véritable signification lorsqu'on fait intervenir la décomposition des idéaux en idéaux primaires; la multiplicité correspondant à un idéal primaire est alors définie comme le nombre de classes modulo cet idéal, linéairement indépendantes, du complément; cela sera traité ailleurs. (E. N.)}, fait qui n'était connu jusqu'à présent que dans le cas spécial où la résultante peut être définie pour des coefficients indéterminés\footnote{Cf. Macaulay, \emph{The algebraic theory of modular systems}, Cambridge Tracts 19 (Cambridge University Press, 1916), no~67; voir aussi Lasker, Zur Theorie der Moduln und Ideale, Math. Ann. 60 (1905), p.~20--116; p.~98. Que, dans ce cas, la résultante et la forme résultante coïncident est montré dans les théorèmes encore non publiés de Hentzelt mentionnés sous 1). (E. N.)}.

Pour déduire ces résultats, on donne aux \S{}~1 et \S{}~2 des théorèmes sur les modules de formes linéaires dont les coefficients sont des nombres rationnels entiers ou des polynômes d'une indéterminée. Le lien avec les idéaux de polynômes introduits au \S{}~3 est assuré par le théorème d'isomorphie du \S{}~4. Le \S{}~5 donne la seconde partie du théorème principal mentionnée ci-dessus, le \S{}~7 le théorème sur les zéros, que précède au \S{}~6 une théorie générale de l'élimination. Là, en même temps, lors de l'introduction de nouvelles indéterminées, est démontrée la décomposition de la résolvante en formes linéaires de ces indéterminées; et ainsi, pour l'élimination donnée ici, est fournie une démonstration irréprochable d'une affirmation de Kronecker\footnote{La tentative de démonstration chez König, \emph{Algebraische Größen} (Leipzig 1903, Teubner), V, \S{}~4 -- pour la théorie de l'élimination de Kronecker -- est notoirement restée sans succès. Que, même pour les multiplicités introduites par König dans la théorie de l'élimination de Kronecker, la seconde partie du théorème principal de Hentzelt ne vaille pas, est montré par Macaulay, loc. cit., p.~23; il y est en outre montré que la théorie de l'élimination de Kronecker ne correspond pas à la décomposition en idéaux primaires. (E. N.)}.

\section*{\S{} 1. Modules de formes linéaires.}

Par grandeurs entières $a,b,c,\ldots$ on entend, dans les deux premiers paragraphes, des nombres rationnels entiers ou des polynômes d'une indéterminée à coefficients dans un corps $P$ quelconque, défini abstraitement; par formes linéaires $a(\xi),b(\xi),\ldots$, de telles formes en un nombre fini d'indéterminées $\xi_1,\ldots,\xi_k$ avec des grandeurs entières pour coefficients.

Un \emph{module $\Amod$ de formes linéaires} est défini par les conditions suivantes: avec $a(\xi)$ et $b(\xi)$, $\Amod$ contient aussi la différence $a(\xi)-b(\xi)$; avec $a(\xi)$, il contient aussi $c\cdot a(\xi)$, où $c$ désigne une grandeur entière arbitraire.

Par le \emph{rang de $\Amod$}, nous entendons, comme d'habitude, le nombre maximal de formes linéaires linéairement indépendantes de $\Amod$; par le $\lambda$-ième \emph{diviseur déterminantiel} $d_\lambda$ de $\Amod$, le plus grand diviseur commun de tous les déterminants de degré $\lambda$ de $\Amod$ (c'est-à-dire de tous les déterminants de degré $\lambda$ que l'on peut former à partir de la matrice des coefficients de $\lambda$ formes linéaires quelconques de $\Amod$). Si $\varrho$ est le rang de $\Amod$, de sorte que $d_1,\ldots,d_\varrho$ sont les diviseurs déterminantiels non nuls, les \emph{diviseurs élémentaires} de $\Amod$ sont définis par $e_1=d_1$, $e_2=d_2/d_1,\ldots,e_\varrho=d_\varrho/d_{\varrho-1}$, $e_{\varrho+i}=0$. On sait que $\Amod$ est un module fini qui possède une base du module formée d'exactement $\varrho$ formes linéaires $a_1(\xi),\ldots,a_\varrho(\xi)$, au moyen desquelles toute forme linéaire de $\Amod$ se représente linéairement, avec des grandeurs entières comme coefficients: $\Amod=(a_1(\xi),\ldots,a_\varrho(\xi))$. Si $A$ désigne la matrice des coefficients de ces formes linéaires, les diviseurs déterminantiels et élémentaires de $\Amod$ coïncident donc avec ceux de $A$; il en résulte d'abord que \emph{chaque diviseur élémentaire $e_i$ divise le suivant $e_{i+1}$}. L'équation matricielle existant d'après la théorie des diviseurs élémentaires,
\[
 P A Q=\begin{pmatrix}
 e_1&&&0\\
 &e_2&&\\
 &&\ddots&\\
 0&&&e_\varrho
 \end{pmatrix}
\]
montre en outre -- lorsque, par une transformation unimodulaire $\xi=Q(\eta)$, on introduit de nouvelles indéterminées $\eta_1,\ldots,\eta_k$ -- l'existence de l'équation de modules:
\begin{equation}
 \Amod=(e_1\eta_1,e_2\eta_2,\ldots,e_\varrho\eta_\varrho).
\tag{1}
\end{equation}

Le concept le plus essentiel pour ce qui suit est celui de module fondamental\footnote{Cf. Steinitz, Rechteckige Systeme und Moduln in algebraischen Zahlkörpern II. Math. Ann. 72 (1912), p.~297--345, no~36. Hentzelt semble en tout cas être parvenu, indépendamment de Steinitz -- avec qui l'on trouve par ailleurs aussi des points de contact --, pour son cas plus simple, à ce concept sous la forme de la formule (4). (E. N.)}, selon la

\medskip
\noindent\textbf{Définition I.} \emph{Un module $\Gmod$ s'appelle module fondamental s'il n'a aucun diviseur propre de même rang}\footnote{Diviseur entendu au sens des modules: $\Amod$ est divisible par $\Bmod$, $\Amod\zero{\Bmod}$, si tout élément de $\Amod$ est contenu dans $\Bmod$; $\Bmod$ s'appelle diviseur propre s'il contient des éléments différents de ceux de $\Amod$.}.

$\Gmod$ est donc un module fondamental si et seulement si, de
\begin{equation}
 c\cdot g(\xi)\zero{\Gmod};\quad c\ne0 \quad\text{on déduit toujours:}\quad g(\xi)\zero{\Gmod}.
\tag{2}
\end{equation}

Le lien entre un module arbitraire et un module fondamental est donné par le

\medskip
\noindent\textbf{Théorème I.} \emph{Tout module $\Amod$ est divisible par un et un seul module fondamental $\Gmod$ de même rang; $\Gmod$ est défini comme le plus petit diviseur de même rang de $\Amod$ et représenté par}
\begin{equation}
 \Gmod=(\eta_1,\ldots,\eta_\varrho);
\tag{3}
\end{equation}
\emph{$\Gmod$ sera appelé le module fondamental de $\Amod$.}

La condition que $\Gmod$ doive être le plus petit diviseur de même rang s'exprime ainsi: $\Gmod$ contient toutes et seulement les formes linéaires $g(\xi)$ qui satisfont à une relation
\begin{equation}
 b\cdot g(\xi)\zero{\Amod};\quad b\ne0.
\tag{4}
\end{equation}
Il en résulte d'abord que $\Gmod$ est un module, puisque, avec
\[
 b_1g_1(\xi)\zero{\Amod};\quad b_1\ne0;
 \qquad
 b_2g_2(\xi)\zero{\Amod};\quad b_2\ne0
\]
on a aussi
\[
 b_1b_2\bigl(g_1(\xi)-g_2(\xi)\bigr)\zero{\Amod};\quad b_1b_2\ne0
\]
et bien entendu aussi $b_1\cdot c g_1(\xi)\zero{\Amod}$. Mais $\Gmod$ est également un module fondamental; car de
\[
 c\cdot g(\xi)\zero{\Gmod};\quad c\ne0
\]
on déduit, d'après (4),
\[
 b c g(\xi)\zero{\Amod};\quad bc\ne0,
\]
donc, de nouveau d'après (4), aussi $g(\xi)\zero{\Gmod}$.

Il n'existe en outre aucun module fondamental $\Gmod_1$ différent du plus petit diviseur de même rang $\Gmod$ qui satisfasse aux conditions. Car $\Gmod_1$ devrait être divisible par $\Gmod$, mais, comme module fondamental, il n'a aucun diviseur propre de même rang, et il est donc identique à $\Gmod$. Enfin, le module représenté par le membre de droite de (3) est fondamental, puisque tout diviseur propre doit contenir une indéterminée supplémentaire $\eta$ et possède donc un rang plus élevé; ainsi (3) donne le module fondamental de $\Amod$ défini de manière unique, ce qui \emph{démontre le Théorème I dans toutes ses parties}.

Un autre concept essentiel pour la suite est le concept dedekindien du quotient de deux modules\footnote{Chez Dedekind, il s'agit de modules de nombres, pour lesquels une multiplication est aussi définie, de sorte que le quotient de Dedekind ne coïncide pas avec le quotient défini ici. (E. N.)}, selon la

\medskip
\noindent\textbf{Définition II.} \emph{Le quotient $\cc=\Amod/\Bmod$ de deux modules de formes linéaires est défini comme l'ensemble de toutes les grandeurs entières $c$ qui satisfont à la condition}
\[
 c\cdot\Bmod\zero{\Amod}
\]
\emph{$\cc$ représente un idéal de grandeurs entières, donc un idéal principal. De même, le quotient $\mathfrak C=\Amod/\bb$ d'un module $\Amod$ de formes linéaires par un idéal $\bb$ de grandeurs entières est défini comme l'ensemble des formes linéaires $c(\xi)$ qui satisfont à la condition}
\[
 c(\xi)\cdot\bb\zero{\Amod}
\]
\emph{$\mathfrak C$ représente de nouveau un module de formes linéaires, et même, dès que $\bb$ est différent de l'idéal nul, un module de même rang que $\Amod$.}

Il faut seulement remarquer qu'il est clair, d'après la définition, que la propriété de module appartient chaque fois au quotient; or les systèmes de grandeurs entières qui possèdent la propriété de module sont des idéaux.

Le concept de quotient conduit à un lien entre module fondamental et plus haut diviseur élémentaire, selon le

\medskip
\noindent\textbf{Théorème II.} \emph{Entre le module fondamental et le plus haut diviseur élémentaire de $\Amod$ subsiste la relation réciproque}
\begin{equation}
 (e_\varrho)=\Amod/\Gmod;\qquad \Gmod=\Amod/(e_\varrho),
\tag{5}
\end{equation}
\emph{où $(e_\varrho)$ désigne l'idéal principal dérivé de $e_\varrho$.}

En effet, puisque chaque diviseur élémentaire divise $e_\varrho$, on obtient, d'après (1) et (3),
\begin{equation}
 e_\varrho\Gmod\zero{\Amod}\quad \text{et par conséquent}\quad (e_\varrho)\cdot\Gmod\zero{\Amod};
\tag{6}
\end{equation}
ainsi $(e_\varrho)$ est divisible par le quotient $\Amod/\Gmod$. Mais la réciproque vaut également; car de
\[
 b\Gmod\zero{\Amod}
 \quad\text{on déduit}\quad b\eta_\varrho\zero{\Amod}
 \quad\text{et donc}\quad b\zero{(e_\varrho)},
\]
ce qui démontre la première formule (5)\footnote{Si l'on pose $\Amod_0=\Amod,\ldots,\Amod_\lambda=(\Amod,\eta_\varrho,\ldots,\eta_{\varrho-\lambda+1})$, de sorte que chaque $\Amod_i$ possède le plus haut diviseur élémentaire $e_{\varrho-i}$, les mêmes raisonnements donnent:
\[
 (e_\varrho)=\Amod_0/\Amod_1,\ldots,(e_{\varrho-\lambda})=\Amod_\lambda/\Amod_{\lambda+1},\ldots,(e_1)=\Amod_{\varrho-1}/\Amod_\varrho.
\]
Plus généralement, soit posé $\zeta_\lambda=b_{\lambda1}\eta_1+\cdots+b_{\lambda\lambda}\eta_\lambda$, soit $(b_{\lambda\lambda},e_\lambda)=1$, et soit
\[
 \Bmod_0=\Amod,\ldots,\Bmod_\lambda=(\Amod,\zeta_\varrho,\ldots,\zeta_{\varrho-\lambda+1});
\]
alors $\Bmod_i$ possède aussi le plus haut diviseur élémentaire $e_{\varrho-\lambda}$, et l'on obtient donc:
\[
 (e_{\varrho-\lambda})=\Bmod_\lambda/\Bmod_{\lambda+1}.
\]}. (6) montre en outre que $\Gmod$ est divisible par le quotient $\Amod/(e_\varrho)$; ce quotient est un diviseur de même rang de $\Amod$, donc, d'après la définition du module fondamental de $\Amod$, il est réciproquement divisible par $\Gmod$, ce qui démontre aussi la seconde formule (5).

Si $d$ désigne un déterminant quelconque de degré $\varrho$ de $\Amod$ qui ne s'annule pas identiquement, alors $d$ est divisible par le $\varrho$-ième diviseur déterminantiel $d_\varrho$, et par conséquent par $e_\varrho$; on a donc $d\Gmod\zero{\Amod}$ et, d'après le raisonnement précédent, il s'ensuit que $\Gmod=\Amod/(d)$. Mais il est essentiel pour la suite que cette dernière représentation de $\Gmod$ comme quotient puisse aussi se démontrer sans revenir aux diviseurs élémentaires, et qu'elle ait par conséquent une validité plus générale, selon le

\medskip
\noindent\textbf{Théorème III.} \emph{Si l'on entend par grandeurs entières des polynômes en plusieurs indéterminées}\footnote{En fait on utilise seulement que les grandeurs entières se reproduisent par addition, soustraction et multiplication, avec validité des règles ordinaires du calcul, c'est-à-dire qu'elles forment un anneau; et de plus que, dans cet anneau, un produit ne s'annule que lorsqu'un facteur s'annule, de sorte que l'anneau peut être étendu à un corps par formation de quotients (adjonction de couples d'éléments). En revanche, aucune représentation du module par une base n'est utilisée; (7) vaut donc par exemple aussi pour le domaine de tous les nombres algébriques entiers comme coefficients. (E. N.)}, \emph{alors la représentation comme quotient}
\begin{equation}
 \Gmod=\Amod/(d),
\tag{7}
\end{equation}
\emph{vaut encore, où $d$ désigne un déterminant quelconque de degré $\varrho$ de $\Amod$ qui ne s'annule pas identiquement.}

Il est d'abord clair que les concepts de rang, de module fondamental et de quotient de modules -- qui maintenant ne devient seulement plus un idéal principal -- se conservent dans cette extension, de même que le Théorème I, à l'exception de la représentation par une base.

Soient maintenant
\[
 a_1(\xi)=a_{11}\xi_1+\cdots+a_{1k}\xi_k,\ldots,
 a_\varrho(\xi)=a_{\varrho1}\xi_1+\cdots+a_{\varrho k}\xi_k
\]
$\varrho$ formes linéaires linéairement indépendantes de $\Amod$, qui en général ne formeront pas une base du module; les indéterminées $\xi$ soient désignées de telle manière que
\[
 d=|a_{ij}|\ne0;\qquad i,j=1,2,\ldots,\varrho.
\]
Il en résulte que $\varrho$ formes linéaires de forme spéciale appartiennent à $\Amod$:
\begin{equation}
 d\xi_\lambda+b_{\lambda,\varrho+1}\xi_{\varrho+1}+\cdots+b_{\lambda,k}\xi_k\zero{\Amod}
 \quad(\lambda=1,2,\ldots,\varrho).
\tag{8}
\end{equation}
Mais $\Amod$ ne peut contenir aucune forme linéaire dépendant seulement de $\xi_{\varrho+1},\ldots,\xi_k$, par exemple $c_{\varrho+1}\xi_{\varrho+1}+\cdots+c_k\xi_k$, car sinon au moins un déterminant à $(\varrho+1)$ rangées, à savoir $d\cdot c_\alpha$, serait différent de zéro.

Or le quotient $\Amod/(d)$, comme diviseur de même rang de $\Amod$, est divisible par $\Gmod$; mais la réciproque vaut aussi. En effet, pour tout $g(\xi)\zero{\Gmod}$, on a par définition:
\[
 c\cdot g(\xi)\zero{\Amod};\quad c\ne0
 \quad\text{et par conséquent}\quad d\cdot c\cdot g(\xi)\zero{\Amod}.
\]
D'après (8), cependant,
\[
 d\cdot g(\xi)=g_{\varrho+1}\xi_{\varrho+1}+\cdots+g_k\xi_k\zero{\Amod},
\]
d'où il vient
\[
 c g_{\varrho+1}\xi_{\varrho+1}+\cdots+c g_k\xi_k\zero{\Amod},
\]
donc, d'après ce qui vient d'être remarqué, $c\cdot g_{\varrho+1}=0,\ldots,c\cdot g_k=0$, et ainsi, puisque $c\ne0$, aussi $g_{\varrho+1}=0,\ldots,g_k=0$, donc $d\cdot g(\xi)\zero{\Amod}$, ce qui démontre (7).

Il faut remarquer que $d\cdot g(\xi)\zero{\Amod}$ résulte aussi directement du fait que les deux modules $(a_1(\xi),\ldots,a_\varrho(\xi))$ et $(g(\xi),a_1(\xi),\ldots,a_\varrho(\xi))$ ont le même rang $\varrho$, de sorte que -- en posant $g(\xi)=c_1\xi_1+\cdots+c_k\xi_k$ -- le déterminant à $(\varrho+1)$ rangées
\[
\begin{vmatrix}
 a_1(\xi)&a_{11}&\cdots&a_{1\varrho}\\
 \vdots&\vdots&&\vdots\\
 a_\varrho(\xi)&a_{\varrho1}&\cdots&a_{\varrho\varrho}\\
 g(\xi)&c_1&\cdots&c_\varrho
\end{vmatrix}
\]
s'annule. Mais la démonstration donnée ci-dessus revient au \S{}~4, à la suite de la formule (30).

\clearpage
\ifdefined\NoetherPaperTwentyTwoSOneBodyOnly
\else
\end{document}
\fi
